% !TeX encoding = UTF-8
\documentclass[institutional,screen,none]{ua-engineering-v6-beamer}
% !TeX encoding = UTF-8
\UASetUniversity{Universidad Austral}
\UASetFaculty{Facultad de Ingeniería}
\UASetDepartment{Departamento de Computación}
\UASetCampus{Pilar}
\UASetCareer{Ingeniería Informática}
\UASetCommission{Comisión A}
\UASetProfessor{Equipo docente}
\UASetSemester{Segundo cuatrimestre 2026}
\UASetCourse{Sistemas Distribuidos}
\UASetDate{2026}


\UASetClassNumber{16}
\UASetClassTitle{Defensa final ante comité}
\title{\UACourse}
\subtitle{Clase 16 --- Defensa final de sistema, objeciones y rúbrica}

\author{\UAProfessor}
\date{\UADate}
\begin{document}

{
\setbeamertemplate{footline}{}
\begin{frame}[plain]
\titlepage
\end{frame}
}

\UASectionSlide{Cierre del curso}

\begin{frame}{Del conocimiento a la defensa}
Durante el curso vimos:
\begin{itemize}
\item mecanismos
\item arquitecturas
\item observabilidad
\item resiliencia
\item confiabilidad
\item incidentes
\end{itemize}

\pause

\begin{block}{Pregunta final}
¿Podés defender integralmente un sistema distribuido ante un comité técnico hostil?
\end{block}
\end{frame}

\begin{frame}{Defender no es describir}
Una defensa final no es:
\begin{itemize}
\item leer slides
\item enumerar tecnologías
\item mostrar un diagrama bonito
\end{itemize}

\pause

Es:
\begin{itemize}
\item formular el problema
\item justificar decisiones
\item reconocer límites
\item responder objeciones
\end{itemize}
\end{frame}

\UASectionSlide{Estructura de la presentación}

\begin{frame}{Secuencia recomendada}
\begin{enumerate}
\item problema y contexto
\item requisitos y restricciones
\item supuestos de falla
\item arquitectura propuesta
\item datos, consistencia y comunicación
\item resiliencia y operación
\item trade-offs
\item límites y trabajo futuro
\end{enumerate}
\end{frame}

\begin{frame}{Marco del curso}
\[
D = (P, F, G, S, T)
\]

\pause

Toda defensa sólida debería hacer explícitos:
\begin{itemize}
\item problema
\item fallas
\item garantías
\item solución
\item trade-offs
\end{itemize}
\end{frame}

\UASectionSlide{Qué mira un comité}

\begin{frame}{Criterios reales}
Un jurado técnico suele mirar:
\begin{itemize}
\item claridad conceptual
\item coherencia interna
\item realismo de supuestos
\item calidad de trade-offs
\item capacidad de responder preguntas
\end{itemize}
\end{frame}

\begin{frame}{Lo que penaliza}
\begin{itemize}
\item slogans sin justificación
\item promesas incompatibles
\item falta de modelo de falla
\item operación ausente
\item respuestas evasivas
\end{itemize}
\end{frame}

\UASectionSlide{Objeciones típicas}

\begin{frame}{Preguntas hostiles frecuentes}
\begin{itemize}
\item ¿por qué esta arquitectura y no otra?
\item ¿por qué no un monolito?
\item ¿qué pasa bajo partición?
\item ¿qué garantía ofrece realmente?
\item ¿cómo opera esto en producción?
\item ¿dónde están los límites?
\end{itemize}
\end{frame}

\begin{frame}{Cómo responder bien}
\begin{itemize}
\item reconocer supuestos
\item explicitar prioridades
\item no defender perfección
\item mostrar consistencia entre dominio y diseño
\end{itemize}

\pause

\begin{block}{Clave}
Una buena respuesta técnica no elimina el trade-off: lo hace visible y lo justifica.
\end{block}
\end{frame}

\UASectionSlide{Revisión crítica}

\begin{frame}{Revisión crítica}
Una defensa madura también puede decir:
\begin{itemize}
\item qué parte es frágil
\item qué parte es más costosa
\item qué parte requeriría validación experimental
\item qué parte dejaría como trabajo futuro
\end{itemize}
\end{frame}

\begin{frame}{No decir “no sé” mal}
Decir:
\begin{quote}
“Eso requiere una decisión adicional de diseño; bajo estos supuestos elegiría X por estas razones”
\end{quote}

es mejor que improvisar una respuesta falsa o totalizante.
\end{frame}

\UASectionSlide{Rúbrica}

\begin{frame}{Rúbrica de evaluación}
\begin{itemize}
\item formulación del problema
\item calidad del diseño
\item tratamiento de fallas
\item trade-offs
\item observabilidad y operación
\item defensa oral
\end{itemize}
\end{frame}

\begin{frame}{Cierre del curso}
\begin{block}{Idea final}
Saber sistemas distribuidos no es repetir mecanismos: es poder defender una arquitectura coherente bajo preguntas difíciles.
\end{block}
\end{frame}


\begin{frame}{Caso de diseño: decidir antes de nombrar tecnología}
\small
\begin{block}{Escenario}Un servicio crítico debe seguir operando ante latencia variable y fallas parciales. El equipo propone ``agregar más infraestructura'' sin declarar qué propiedad necesita preservar.\end{block}
\begin{enumerate}
\item ¿Cuál es el invariante o garantía que realmente importa?
\item ¿Qué falla concreta amenaza esa garantía?
\item ¿Qué mecanismo de Defensa final ante comité es pertinente y cuál sería sobreingeniería?
\item ¿Qué métrica permitiría comprobar el costo de la decisión?
\end{enumerate}
\end{frame}
\begin{frame}{Checkpoint de comprensión}
\small
Al terminar esta clase deberías poder:
\begin{itemize}
\item explicar el problema sin apoyarte en nombres de productos;
\item identificar el supuesto de falla más importante;
\item distinguir una propiedad de seguridad de una de progreso;
\item construir un escenario donde la solución deja de ser suficiente;
\item defender el trade-off ante una objeción técnica.
\end{itemize}
\end{frame}

\begin{frame}{Simulación de comité}
Cada equipo recibe una objeción de consistencia, una de falla y una de operación. Tiene 60 segundos para declarar supuesto, decisión y trade-off antes de profundizar.
\end{frame}
\begin{frame}{Evidencia por encima de opinión}
Una defensa excelente enlaza cada afirmación importante con una traza, prueba, benchmark, fault injection, ADR o resultado reproducible.
\end{frame}
\begin{frame}{Cierre formativo}
El curso culmina cuando el estudiante puede reconocer el tipo de problema, seleccionar un mecanismo proporcional, medir su comportamiento y revisar su decisión ante nueva evidencia.
\end{frame}

\end{document}
